\documentclass[12pt]{article}  
\newcommand{\say}[1]{``#1''}  
\newcommand{\nsay}[1]{`#1'}  
\usepackage{endnotes}  
\newcommand{\B}{\backslash{}}  
\renewcommand{\,}{\textsuperscript{,}}  
\usepackage{setspace}   
\usepackage{tipa}  
\usepackage{hyperref}  
\begin{document}  
\doublespacing  
\section{\href{tapestery-embroidery.html}{On Tapestry Embroidery}}  
First Published: 2026 June 19

\section{Draft 4: 19 June 2026}

I'm working on a design for a tapestry embroidery.  
My plan is to have some sort of a color gradient beneath a loose woven layer beneath some sort of geometric design.  
Part of me would like to play with transparency between layers, but that feels better as a follow-up project once I've demonstrated to myself that I can do something like this at all.

There are a series of questions I need to answer:

\begin{enumerate}  
\item Material:

\begin{enumerate}

\item What count Aida fabric?

Right now I think that I have a large roll of 18 squares per inch, so that's a tempting choice. However, I think that the last project I did was on 14 count Aida?

\item How many threads per stitch?

In the previous project, I used 6. In some quick trials, I think I found that four looked best to me. Will need to see, though.

\item How many colors?

DMC makes an almost infinite amount of colors. The more I use, the more smooth of a gradient I can produce, but also the more pain it will be, because I'll need to swap colors more. Also, there's something to be said for the digital nature of stitch-work.

\item How large of a piece of fabric do I need?

This question is going to need to be answered based on the first Material question and the first design question. Ideally, smaller is usually better, but there's also something beautiful in a large embroidery.

\end{enumerate}

\item Design:

\begin{enumerate}

\item How many stitches large?

\item What base layer color or color gradient?

\item Weave:

\begin{enumerate}

\item How many stitches thick are the strands?

\item What's the strand density? I.e. how much of the base gradient is visible between crossings?

\item What color or colors are the weave strands?\footnote{look at that, managed to write the sentence so I didn't need an is/are}

\item Is the weave going to be a Celtic knot? If so, what is the design for it?

\item Do we incorporate the base gradient colors in the weave strands, and if so, how much?\footnote{initially this said \say{do we need to} but I realize that this project is not needed in any sense, and so there is no version wherein it is necessary to have it}

\end{enumerate}

Of course, all of these questions are fairly inter-reliant.  
A thicker strand will result in a higher strand density, and a larger total stitch-count project will likely mean that I need to have a thicker strand for visibility.  
Then again, there's also something really compelling to me in the idea that the strands might be almost invisible when looking at the embroidery from afar.  
Who can say for certain?

\item Foreground:

\begin{enumerate}

\item Do I want braids?

\item Do I want a Celtic Knot?

\item What shape?

\item What thickness?

\item How gappy?\footnote{I'm trying to think of a word that's like saturation but pixel-wise. E.g. if I have two lines 10 cm apart from center to center, the gap will be smaller if the lines are thicker. What's the term for the relative size of the gap? Eh. Gappy with this footnote works}

\item What colors? I think that it could be cool to do something that gives an impression of depth, so like crossing gets darker.

\end{enumerate}

And, of course, I want the top layer to appear organic and not rigid, even if that's going to be really difficult to do.

\end{enumerate}

\end{enumerate}

There are more than a few trials that I'll need to do as well to figure out the answer to all of these questions.  
For now, though, I think that I'm going to work on trials for the base weave question.

That is, how thick and how dense, and then by extrapolation, how much that affects the visibility of the background.  
I have a vague idea, if nothing else, of how to make the weave design, so I need to think of the trials I'm going to run.  
For now, I'm thinking I just pick a relatively bold color weave against a contrasting background color.  
Then try one through four stitches thick per strand, with the same number of stitches per strand as the gap between strands, making four or so crossings?  
That should take up some time, so let's go ahead and draw it out.

Oh, shoot, part of this will also end up depending on the number of stitches per inch that I choose. I'm just going to commit to the fabric I have with me being the fabric count I'll be using for the final project?  
Eh, the relative density probably doesn't depend too terribly much on the thread fabric saturation.  
Let's just not worry about it for now.

Woo! Blog post ready to post.

\section{Draft 3: 19 June 2026}

I don't think that I wrote anything about my design process for the last embroidery project I did.  
This one, however, I will document here, at least right now.  
I make no promises about continuing the documentation, but I hope that my writings here might be useful to a future me or other who would like to also make a design for embroidery.\footnote{shoot, the Buddhists are gonna be coming after me for that line. There is no static me, so the future me is an other}

The planned design is a partially transparent flame and logs atop a woven pattern atop a sky gradient.

Open design questions:

\begin{enumerate}

\item How many total colors will I use?  
\item How large will this be in stitches?  
\item What size fabric am I using? Probably 18 holes per inch Aida, but maybe not?  
\item How many threads per stitch. If using 18, then the answer is four, I think. If using 14, I think that the answer becomes six? Will need to run some tests to see what does what I need.  
\item What gradient, exactly? I think that a sunrise to sunset motif would be beautiful, but need to figure out what exactly that will look like. Potential options include:  
\begin{itemize}

\item Center of fabric white, corners black. Diagonal corners for sunrise (warm colors) and sunset (cold colors) both into the blue of the morning and then into the white of the noonday sun  
\item Center of fabric black, corners white, the same but in reverse.  
\item Single gradient across the fabric from black to blue to black or vice versa for sunrise sunset  
\item Rainbow  
\item Something more chaotic?  
\item Seasons? So like green into red into white?  
\item Solid, monochrome color  
\end{itemize}

Right now I'm mostly leaning towards something that goes a variant of dark to light via warm colors back to dark via cold or have bright to dark to bright.

\item What weave density to use? I.e. how much of the gradient is visible in the gaps between weave lines?  
\item What weave thickness? I.e. how many stitches per weave line?  
\item Pure weave or make it a Celtic knot? If knotted, what knot design?  
\item What color for the weave?  
\item Does the weave line need to incorporate the gradient within it? E.g. if the weave is silver and it should be on top of black, do I do like two silver and two black?  
\begin{enumerate}  
\item If it does, what percentage?  
\end{enumerate}

\item How do the fire and logs look?  
\item Smoke?  
\item Etc.

\end{enumerate}

So, things to trial out now:

\begin{enumerate}

\item Weave densities.  
\item Weave thicknesses.  
\item Weave and background color ratios within the weave thread

\end{enumerate}

These trials are of course all interlinked variables. The thickness is particularly going to depend on the final design as well.  
I'll need to decide how visible I want the weave to be and how thick relative to the total size of the fabric I want it to be.

You know what, I don't actually think that I want the art to be representational.  
I just want geometry.  
So, then, I'll do a cute flowy Celtic knot made of a braid! That will be gorgeous.

Draft four for new set of questions.

\section{Draft 2: 19 June 2026}  
I apologize for my absence.  
Anyways.

I'm wanting to make a tapestry embroidery project.  
My plan is to have some sort of color gradient as the base layer, a weaving or potentially even Celtic knotwork secondary layer, and then a flowing Celtic knot made of a braid as the top layer.  
My other idea was to have something like partially transparent glasses, and then attempt to make something like \say{the view of the floor from a lit low-saturation stained glass window}, but I think that might be a good third or fourth project.

So, what is keeping me from making the pattern?

First, I need to decide what color gradient to have as a base layer.  
Right now I'm drawn to the idea of black at the corners and white in the middle\footnote{or vice versa, depending} with two sides being sunrise color gradients\footnote{read: warm colors of red and yellow} and the other two being sunset\footnote{read: purples and maybe greens? Because I like green, at least} transitioning into the blue of a clear day and then the whiteness of staring into the sun.

The second question is how to do the weave\footnote{even if it ends up being Celtic knots, I will be having it as a fairly rigid pattern, and so I'm comfortable labeling it as the weave for the purposes of this project}.  
Depending on the density of the weave, it may be vitally important that I have the base color represented in it.  
Questions within this question include whether I can stitch multiple times and if that behaves the same as stitching once with multiple threads.  
E.g. I might want one strand of silver. If I just stitch the entire canvas in silver and then use three more threads for the appropriate colors, will that look the same as just swapping colors? I guess not, and also it feels icky to do that.

Third question is what the upper layer will be.

I usually feel inspired by flames and fire, and so kind of want that.  
I also love Celtic design, and want to incorporate some of that?  
Or, at least, I love patterns, and so that feels like a good thing to do.  
As a Catholic of Irish descent, I feel no qualms about using these motifs in my work.  
While I do love mandalas, or at least the Westernized versions I see in adult coloring books.\footnote{looking it up on Wikipedia, very few of the shown examples resemble those in any real sense}

However, as I mention in the first draft, I am also interested in potentially exploring some sort of representational piece of art.  
Fire would work well for that, as would a flame with smoke.  
If I did a fire, though, there's a potential issue with the background gradient being lost?  
Though, fire is often at least somewhat transparent.

Hmm.  
Hard to say for certain.

Now, one benefit of doing all of this pre-work is that the number of stitches I will need to do is completely independent of the complexity of the design, which is a huge perk to me.  
Like with cross stitch, each stitch is effectively a pixel.  
Since I'm planning to use every pixel, the only question is how compressible the pixels will be.\footnote{in a data compression sense. The more similar the colors are next to each other, the more easily the computer can just go \say{that whole area is one color, just fill in X,Y box}}  
And, I enjoy this part of the planning too.

Another question I have is how many colors to use.  
The last embroidery I did was with non-DMC colors\footnote{spicy}, because I bought a sunburst kit from the store.  
That had a limited number of colors, which meant I had to and got to be creative about my choices.  
I did a six-strand embroidery, but doubled up each strand\footnote{thanks for the tip friend!} so that I wouldn't have to worry about the thread falling off the end of the line.  
That meant I had up to three color options, which did make it really pretty.

In general, I don't see a lot of people mixing and matching different colors of thread, which I'm not totally sure I understand.  
Sure, Yellow and off Yellow may approximate Yellow three, but the interplay of which exact thread is on top for any given part of the embroidery is unknown until stitching, which gives it more animation.  
Then again, if the goal is to draw attention away from the thread, I guess I can see why it is that people would use the specific color they want.

Oh! I remember what I don't like about Bargello now!

Because it's worked vertical, the Aida fabric vertical strips often also shines through.  
Tapestry embroidery, by contrast, is worked diagonal, so is able to cover the full thing.

Not for this project, but for a future project, I'm also very curious about layering tapestry embroidery.  
That is, if I embroider a base layer all going like a forward slash, if I then use a single strand of floss to go in the back slash direction, will I be able to see two images?  
I don't know, and I am more than a little curious.

Back to the point, though: what do I want to have as the major image of the picture?

It needs to be something with gaps so that the gradient can shine through.  
It needs to be representational, because I decided it does.  
It needs to be pretty.  
  
I guess I could do four layers and in such a way still have some sort of a Celtic triquetra as the third layer?  
Eh.

Maybe text as the top layer?  
That could be really nice, actually.  
What would the text say?

Hmm.  
There are so many beautiful texts in the world.

Maybe just a classic \say{Love}, or something like that?  
I could make the most fun commentary on modern art with three tapestry embroideries that are super detailed that in total say \say{Live, Laugh, Love}, but that feels maybe a little too on the nose for my taste.

A net?  
That's got some promise, if I do like a net that's been draped over something.

Something burning?

A tree branch on fire?

A lightning strike?

I think that a partially transparent flame does actually seem like my favorite idea.  
In the next draft, I will write down the series of questions that I actually need to answer here.

\section{Draft 1: 19 June 2026}

So, I haven't been here in a long while.  
I've even written a few posts since the last time I was here, but none felt up to the quality of a returning post.  
However, I'm currently disabled\footnote{officially. In practice, I stand by a social model of disability which tells me that, since I have the accommodations that I need right now, I'm not disabled. Legally and practically, though, I cannot do the things that are required for me to perform labor I have been contracted for, and so am disabled in that sense absolutely.}, and have therefore decided that now is as good a time as any to get back into this site.  
I know that I feel better when I write, and I know that I have multiple\footnote{maybe even more than I can count on one (functional) hand!} readers, and they have expressed sadness that the site lacks updates.

I'm going to go back to calling it a blog and saying that I'm blogging for now, because, while I am interested in the meaning of what I'm doing, I'm also not willing to stop myself and remember what I decided to call the site.  
I am still interested in doing meta-data and whatnot for the site, but that's something that would need to happen in the future.  
I have ideas for how I don't have to use the new blog\footnote{\href{j.rebelsky.com/newblog}{j.rebelsky.com/newblog}} and still write in markdown, but I'm also out of my love of markdown era.  
Ideally I think that there would be a way for me to use mediawiki as the base writing, because I like how it lets me do nested lists the best.  
However, that's not the purpose of this post.

As I assume everyone who reads this blog knows\footnote{if not, hi! I presumably miss you dearly and would love to have you in my life again. Please drop me a line and I'm more than happy to continue the thread!}, I've begun seriously dating someone.  
Among her infinite wondrous qualities, she also has an aunt who was\footnote{is? love living in capitalism end-staging where we are only what labor we perform for the job we are paid for is what defines what we are} a medievalist and does tapestry.  
When I first heard this, I assumed tapestry in the woven on loom sense, which was wondrous and fascinating to me.\footnote{semi-related, I have an idea for a future blog post (which is returning to the end of these docs) about wonder}

I later learned that she does tapestry embroidery.\footnote{I still don't know if she also does loom tapestry, but}  
I loved my Bargello embroidery project, but I had some dissatisfaction with it, much of which I think came down to the vertical stripes, for reasons that I cannot fully articulate here.  
Tapestry embroidery, by contrast, is performed with diagonal stitches, like the first half of a cross stitch.  
As such, it is more something that I cannot find the word for right now.

And so, I decided that I really want to make a tapestry embroidery project.

I recently watched a documentary about Jerry's Map, which is a performance art piece\footnote{maybe?} or an algorithmic art piece\footnote{maybe?} or something.  
It's a series of panels that an artist has been working on for decades.  
It never ends, and part of me is inspired by that same idea, even if I would not want to make a tapestry in the way he does it.\footnote{I don't like the idea of destroying old parts of the art, and I like pre-planning}  
Anyways, that was something I have been thinking about as I think about how to make this tapestry.

Why come back to embroidery now?  
A few reasons.

\begin{enumerate}

\item I know that it's good for me to do creative things  
\item Right now I can't use my right hand for much, and embroidery feels like a good way to get some dexterity in my left hand.\footnote{yes, I do recognize the humor in the fact that dexterity explicitly refers to things being like the right hand. Anyways}  
\item I like making art  
\item I want my internal association with embroidery not to only be the grief embroidery I did  
\item I like showing off the things that I create  
\item I like seeing beautiful things and want to make them

\end{enumerate}

What's stopping me from making the embroidery right now, then?

A few things.\\ Chief among them, however, is the fact that I don't entirely know what I want to make.  
The tapestry that my partner's aunt\footnote{gosh I hope this doesn't count as doxxing. If either of the affected parties here want me to take this down, please let me know} made was a portrait equivalent\footnote{I do know abstract versus non-representational. I don't really know the art words for most of the other things, like the difference between a picture of a person and a portrait, etc}, and that does lead me to think that it might be best to make something representational.  
In my last project, I made a large Celtic knot with a cross in the white space.  
There will not be white space in this project, because tapestry embroidery fills the entire page.\footnote{interestingly, looking up Bargello embroidery, I'm told that it's mathematical in nature, so there's I guess an argument to be made that what I'm making would still be bargello}

I loved the color gradient in the initial embroidery I made, and would like to really lean into that.  
Many parts of my brain love Celtic knot and adjacent patterning, and so I was thinking it could be really nice to have a Celtic knot background\footnote{or at least a diagonal weaving with over and under alternating} underneath whatever I make.  
I don't love that the last project I did was so blocky.  
That's not the right word.

Digital?  
Maybe.

I guess that I'm trying to express that I don't like how rigid\footnote{There it is!} the previous one was.

I was considering perhaps doing a braid or a flowy Celtic Knot\footnote{Like a triquetra, as it's apparently called}, so that I can use something there.  
I think that might be the best bet, so will be doing that!  
Now I just need to do the design portion.

Alrighty, on to draft two, which is hopefully going to be more coherent.

\section{Upcoming posts}  
\begin{itemize}

\item Wonder

\item Curse of knowledge, esp re. Music Theory\footnote{stolen from a prior list}

\end{itemize}

\end{document}